\chapter{Modelling archetypes and their closed forms}

\begin{orient}
\textbf{What this chapter is for.} A modelling problem is two problems stacked on top of each other, and only the second one is calculus. The first is translation: reading a paragraph of English about coffee, brine, capacitors or falling bodies and producing a differential equation with the right variables, the right signs and the right units. The second is solving that equation, and by this point in the book you can already solve every equation the first stage will hand you, because there are only about six of them. Almost all the marks in this subject are lost in the translation, and almost all of the translation is one instruction: name the quantity that is accumulating, then write \emph{rate of change equals rate in minus rate out} and refuse to write anything else until both terms are expressed in the named quantity. This chapter works through the archetypes one by one. For each you get the physical statement, the equation it becomes, the closed form, and the phrase in the problem's wording that tells you which archetype you are looking at. By the end you should be able to read a modelling problem and write down its solution family before you have finished reading the second sentence, and to know which of the standard sign errors you are at risk of committing.
\end{orient}

\section{The translation step}

Every model in this chapter is produced by the same four decisions, made in this order, and it is worth making them explicitly the first fifty times.

\emph{One: name the state variable and its units.} Not ``temperature'' but $T(t)$ in degrees Celsius; not ``salt'' but $A(t)$ in kilograms of dissolved salt, as distinct from $c(t)$ in kilograms per litre. Half of all mixing errors are a confusion between an amount and a concentration, and naming the units prevents it.

\emph{Two: write the balance law.} For anything that accumulates, $\dv{X}{t}=(\text{in})-(\text{out})$. For anything mechanical, $F=ma$ with every force listed and signed against a stated positive direction. For a circuit, Kirchhoff's voltage law around the loop.

\emph{Three: express every term in the state variable and $t$ alone.} This is the step that produces the equation. An outflow term that still mentions concentration is not finished; a drag term that still mentions the distance fallen is not finished.

\emph{Four: read off the type.} Separable, first-order linear, or a linear system. There is no fourth option anywhere in this chapter.

\begin{keynote}[Units are a free proof]
Every constant that appears in a model carries units, and checking them catches structural errors that no amount of rereading will. In $T'=-k(T-T_a)$ the left side is degrees per minute and $T-T_a$ is degrees, so $k$ is per minute and $1/k$ is a time. In $L I'+RI=E$ every term is a voltage, so $L/R$ has units of time and $\tau=L/R$ is forced; nobody who checks this ever writes $\tau=R/L$. In $\dv{A}{t}=r c_{\text{in}}-\frac{r_{\text{out}}}{V}A$ the second term is (litres per minute) times (kilograms per litre), which is kilograms per minute, matching the left side; the wrong version $r_{\text{out}}A$ would be kilogram-litres per minute and is dead on arrival. Do this check before you integrate, not after.
\end{keynote}

The first three archetypes all reduce to the same equation, $y'=\lambda y$ or a shifted version of it, and they are grouped here for that reason. What distinguishes them in practice is not the mathematics but which quantity turns out to be the one decaying exponentially.

\section{The exponential family}

The equation $y'=\lambda y$ is the entire content of the next three techniques. What varies is the substitution that produces it. In decay it is already there. In cooling you must subtract the ambient temperature first. In logistic growth you must form the odds. Each of those is a one-line change of variable, and recognising which one to make is the whole recognition problem.

\tech{Radioactive decay, half-life, and dating}{easy}
\cue The phrases ``half-life'', ``decays at a rate proportional to the amount present'', ``carbon dating'', or any quantity described by a fixed percentage loss per fixed time.
\method $N'=-\lambda N$ gives $N=N_0\eu^{-\lambda t}$, with $t_{1/2}=\dfrac{\ln2}{\lambda}$. In practice do not compute $\lambda$ at all. Write
\[N(t)=N_0\,2^{-t/t_{1/2}},\]
which is the same function with the data already substituted, and answer every question directly from it. Inverting, a sample at fraction $f$ of its original content has age $t=t_{1/2}\log_2(1/f)$.

\begin{worked}
A sample retains $25\%$ of its original carbon-14, whose half-life is $5730$ years. Since $0.25=2^{-2}$, the sample is exactly two half-lives old: $t=11460$ years, with no logarithm required. If instead the sample retains $62\%$, the same formula gives
\[t=5730\cdot\frac{\ln(1/0.62)}{\ln2}=5730\times0.689661=3951.8\ \text{years},\]
so roughly $3950$ years, which is as much precision as a radiocarbon measurement supports.
\end{worked}

\begin{worked}[Decay against a source]
A short-lived isotope is produced in a reactor at a constant rate $R$ atoms per second and decays with constant $\lambda$. Now $N'=R-\lambda N$, first-order linear, and the substitution $U=N-R/\lambda$ removes the constant: $U'=-\lambda U$, so
\[N(t)=\frac{R}{\lambda}+\Bigl(N_0-\frac{R}{\lambda}\Bigr)\eu^{-\lambda t}.\]
The equilibrium $N_\infty=R/\lambda$ is called \emph{secular equilibrium}, and the approach to it has exactly the half-life of the isotope: after one half-life the gap to equilibrium has halved. Starting from $N_0=0$, the inventory reaches $90\%$ of equilibrium after $\log_2(10)\approx3.32$ half-lives.
\end{worked}

\begin{trap}
Reporting the decay constant when the half-life was asked for, or the reverse. They differ by the factor $\ln2\approx0.693$, so both answers are the same order of magnitude and neither looks wrong. Attach units the moment you compute either one: $\lambda$ is per year, $t_{1/2}$ is years, and $\lambda t_{1/2}=\ln 2$ is the relation to quote.
\end{trap}

\begin{seealsob}Newton's law of cooling and the equal-time ratio; logistic growth and the odds substitution; compartment models.\end{seealsob}

Cooling is decay with an offset, and the offset is where the difficulty lives. A cooling problem that tells you the ambient temperature is a decay problem in disguise and takes one line. A cooling problem that withholds the ambient temperature looks impossible and is not, because three readings at equal spacing determine it algebraically.

\tech{Newton's law of cooling and the equal-time ratio}{medium}
\cue A body at temperature $T$ in surroundings at $T_a$, with two or three temperature readings at stated times. Also: ``how long until it reaches'', ``what was the temperature before the measurements began'', ``estimate the room temperature''.
\method The law is $T'=-k(T-T_a)$. Do not solve it in $T$. Set $U=T-T_a$, the \emph{excess}, and note $U'=-kU$, so
\[U(t)=U_0\eu^{-kt},\qquad T(t)=T_a+(T_0-T_a)\eu^{-kt}.\]
The operational consequence is that over any two intervals of equal length the excess is multiplied by the same factor $\rho=\eu^{-k\Delta}$. So if readings are taken at equally spaced times, the excesses form a geometric progression, and $k$ never has to be computed. If $T_a$ is unknown, three equally spaced readings $T_0,T_1,T_2$ determine it: the condition $(T_1-T_a)^2=(T_0-T_a)(T_2-T_a)$ expands to a linear equation, whose solution is
\[T_a=\frac{T_0T_2-T_1^{2}}{T_0+T_2-2T_1}.\]

\begin{worked}[Recovering the room temperature]
A body is measured at $80^{\circ}$, then $60^{\circ}$ ten minutes later, then $48^{\circ}$ after a further ten minutes. The ambient temperature is not given. Apply the formula:
\[T_a=\frac{80\cdot48-60^{2}}{80+48-120}=\frac{3840-3600}{8}=30^{\circ}.\]
Check it against the geometric progression: the excesses are $50$, $30$, $18$, and $30/50=18/30=\tfrac35$, so $\rho=\tfrac35$ per ten minutes. Everything else now follows without solving for $k$. Ten minutes later the excess is $18\cdot\tfrac35=10.8$, so $T(30)=40.8^{\circ}$. And the time to reach $35^{\circ}$, an excess of $5$, satisfies $50\rho^{t/10}=5$, so
\[\frac{t}{10}=\frac{\ln10}{\ln(5/3)}=\frac{2.302585}{0.510826}=4.5075,\qquad t\approx45.1\ \text{minutes}.\]
\end{worked}

\begin{worked}[The two-reading version]
Coffee at $90^{\circ}$ in a $20^{\circ}$ room is $70^{\circ}$ after ten minutes. The excess went from $70$ to $50$, so $\rho=\tfrac57$ per ten minutes, and after twenty minutes the excess is $70\cdot\bigl(\tfrac57\bigr)^{2}=\tfrac{250}{7}$. Hence
\[T(20)=20+\frac{250}{7}=\frac{390}{7}\approx55.71^{\circ}.\]
Note what was not done: $k=\tfrac{1}{10}\ln\tfrac75\approx0.033647$ was never needed, and computing it would have introduced a rounding error into an exact answer.
\end{worked}

\begin{trap}
Applying the ratio to $T$ rather than to $T-T_a$. In the coffee example, $70/90$ is not the decay factor and $90\cdot(7/9)^{2}\approx54.4$ is not the answer, though it is close enough to look right. Temperatures do not decay geometrically; excess temperatures do. The tell is that the wrong method predicts $T\to0$, which is a statement about absolute zero that no coffee cup has ever made.
\end{trap}

\begin{seealsob}Radioactive decay and half-life; RC and RL circuits; compartment models.\end{seealsob}

Unbounded exponential growth is a model with a short shelf life, since every real population runs into a resource ceiling. The standard repair multiplies the growth rate by a factor that vanishes at the ceiling, and the resulting equation is nonlinear but still solvable in closed form, because a single substitution linearises it.

\tech{Logistic growth and the odds substitution}{medium}
\cue $P'=kP(K-P)$ or $P'=rP\bigl(1-\tfrac PK\bigr)$: growth proportional to the population \emph{and} to the room remaining. In words: ``grows exponentially at first but levels off at a carrying capacity''.
\method The two forms are the same equation with $r=kK$. Partial fractions,
\[\frac{1}{P(K-P)}=\frac1K\left(\frac1P+\frac1{K-P}\right),\]
show that the correct variable is the \emph{odds} $W=\dfrac{P}{K-P}$, which satisfies $W'=rW$ exactly. So the odds grow purely exponentially, $W=W_0\eu^{rt}$, and unwinding gives
\[P(t)=\frac{KP_0}{P_0+(K-P_0)\eu^{-rt}}=\frac{K}{1+W_0^{-1}\eu^{-rt}}.\]
Work in odds throughout: every data point converts to an odds value, every question converts back at the end, and no logarithm is ever taken of a difference.

\begin{worked}
A deer population has carrying capacity $K=500$, with $N(0)=100$ and $N(4)=200$. When does it reach $350$?

Convert the data to odds: $W_0=\tfrac{100}{400}=\tfrac14$ and $W(4)=\tfrac{200}{300}=\tfrac23$. The odds multiplied by $\tfrac{2/3}{1/4}=\tfrac83$ over four years, so $\eu^{4r}=\tfrac83$ and $r=\tfrac14\ln\tfrac83=0.245207$ per year. The target $N=350$ is odds $\tfrac{350}{150}=\tfrac73$, which is $\tfrac{7/3}{1/4}=\tfrac{28}{3}$ times the initial odds. Hence
\[t=\frac{\ln(28/3)}{r}=4\,\frac{\ln(28/3)}{\ln(8/3)}=4\times\frac{2.233592}{0.980829}=9.109\ \text{years}.\]
A fourth-order Runge-Kutta integration of $N'=rN(1-N/500)$ from $N(0)=100$ returns $N(9.108995)=350.000$, confirming both $r$ and $t$.
\end{worked}

\begin{worked}[The inflection point is the fastest growth]
Differentiating $P'=rP(1-P/K)$ gives $P''=r P'(1-2P/K)$, which vanishes at $P=K/2$. So the population curve has its steepest point exactly at half capacity, with maximum growth rate $rK/4$. This is worth remembering as a data-fitting shortcut: if a problem states the year of fastest growth, it has told you the year at which $P=K/2$, hence that $W=1$, hence that $\eu^{rt}=W_0^{-1}$ at that time. In the deer example $W=1$ at $t=\ln4/r=5.653$ years, when the herd passes $250$.
\end{worked}

\begin{trap}
Calibrating the parameters in the wrong order, or believing that two data points determine all three of $P_0$, $r$ and $K$. They do not. With $K$ given, two readings fix $W_0$ and $r$ and the problem is closed; with $K$ unknown you need a third condition, and the cleanest one is the inflection time. The second trap is sign: writing $P'=rP(1-K/P)$ instead of $1-P/K$ inverts the whole model and predicts extinction.
\end{trap}

\begin{seealsob}Radioactive decay and half-life; Bernoulli equations; compartment models.\end{seealsob}

\section{Bookkeeping models: rate in minus rate out}

The exponential family had a single reservoir with no explicit flows. The models in this section make the flows explicit, and the extra difficulty is entirely one of accounting: what enters carries a concentration you are told, and what leaves carries the concentration currently inside, which is a function of the unknown. The figure below is the whole subject.

\begin{center}
\begin{tikzpicture}[node distance=7mm and 11mm]
\node[dnode] (in) {Inflow\\ $r_{\text{in}}$ L/min at\\ $c_{\text{in}}$ kg/L};
\node[dtest, right=13mm of in] (t1) {Tank 1\\ $A_1(t)$ kg in $V_1(t)$ L\\ concentration $A_1/V_1$};
\node[dtest, right=13mm of t1] (t2) {Tank 2\\ $A_2(t)$ kg in $V_2(t)$ L\\ concentration $A_2/V_2$};
\node[dnode, right=13mm of t2] (out) {Outflow\\ $r_{\text{out}}$ L/min at\\ $A_2/V_2$ kg/L};
\node[dleaf, below=9mm of t1] (b1) {$A_1'=r_{\text{in}}c_{\text{in}}-\dfrac{r_{12}}{V_1}A_1+\dfrac{r_{21}}{V_2}A_2$};
\node[dleaf, below=9mm of t2] (b2) {$A_2'=\dfrac{r_{12}}{V_1}A_1-\dfrac{r_{21}+r_{\text{out}}}{V_2}A_2$};
\draw[dedge] (in) -- (t1);
\draw[dedge] (t1) -- node[dlab,above]{$r_{12}$} (t2);
\draw[dedge] (t2) -- node[dlab,above]{$r_{\text{out}}$} (out);
\draw[dedge] (t1) -- (b1);
\draw[dedge] (t2) -- (b2);
\end{tikzpicture}
\end{center}

Delete Tank 2 and you have the single-tank problem; keep it and you have a compartment model. Nothing else changes. The volumes are bookkeeping too: $V_i(t)=V_i(0)+(\text{flow in}-\text{flow out})t$, and a tank whose volume changes has a genuinely non-constant coefficient in front of $A_i$, which is the difference between a five-line problem and a fifteen-line one.

\tech{Mixing and tank problems}{core}
\cue Brine, pollutant, or dye entering a well-stirred tank at a stated rate and concentration, and leaving at a stated rate. The words ``well-stirred'' or ``thoroughly mixed'' are there to license the assumption that the outflow concentration equals the tank average.
\method Track the \emph{amount} $A(t)$, not the concentration. Then
\[\dv{A}{t}=r_{\text{in}}c_{\text{in}}-\frac{r_{\text{out}}}{V(t)}A,\qquad V(t)=V_0+(r_{\text{in}}-r_{\text{out}})t.\]
This is first-order linear, always. If $r_{\text{in}}=r_{\text{out}}$ the coefficient is constant and the answer is a steady state plus a decaying transient. If they differ, the integrating factor is a power of $V(t)$: with $r_{\text{in}}-r_{\text{out}}=\delta\ne0$ the coefficient is $\frac{r_{\text{out}}}{V_0+\delta t}$ and the factor is $(V_0+\delta t)^{r_{\text{out}}/\delta}$.

\begin{worked}[Constant volume]
A $200$ L tank holds pure water. Brine at $0.4$ kg/L enters at $5$ L/min and the mixture leaves at $5$ L/min. Here $V\equiv200$, so
\[A'=2-\frac{A}{40},\qquad A(0)=0.\]
The steady state is $A_\infty=80$ kg, which is the tank at the inflow concentration, and the transient decays with time constant $40$ minutes:
\[A(t)=80\bigl(1-\eu^{-t/40}\bigr).\]
Note the two sanity checks: the limiting concentration $80/200=0.4$ kg/L must equal $c_{\text{in}}$, and the time constant $V/r=200/5=40$ minutes is the time to replace the tank's volume once.
\end{worked}

\begin{worked}[Changing volume]
A $100$ L tank holds pure water. Brine at $1$ kg/L enters at $6$ L/min; the well-stirred mixture is drawn off at $4$ L/min. Find the salt content when the volume reaches $200$ L.

The volume grows at $2$ L/min, so $V(t)=100+2t$ and the tank reaches $200$ L at $t=50$ min. The balance is
\[\dv{A}{t}=6-\frac{4A}{100+2t}=6-\frac{2A}{50+t},\]
so with $\mu=(50+t)^{2}$,
\[\bigl[(50+t)^{2}A\bigr]'=6(50+t)^{2}\ \Longrightarrow\ (50+t)^{2}A=2(50+t)^{3}+C.\]
The condition $A(0)=0$ gives $C=-2\cdot50^{3}=-250000$, hence
\[A(t)=2(50+t)-\frac{250000}{(50+t)^{2}},\qquad A(50)=200-\frac{250000}{10000}=175\ \text{kg}.\]
Runge-Kutta integration of the original equation returns $A(50)=175.000$. The concentration at that moment is $175/200=0.875$ kg/L, still short of the inflow value $1$, as it must be.
\end{worked}

\begin{trap}
Writing the outflow term as $r_{\text{out}}c_{\text{in}}$, or as $r_{\text{out}}A$. What leaves carries the tank's \emph{current} concentration $A/V(t)$, so the term is $r_{\text{out}}A/V(t)$; the first error makes the equation trivially separable and wrong, the second is dimensionally impossible. The third and subtlest error is holding $V$ fixed when the rates differ: that turns a variable-coefficient equation into a constant-coefficient one and produces a solution that approaches a finite limit when the true one grows without bound.
\end{trap}

\begin{seealsob}Compartment models; first-order linear equations and the integrating factor; Torricelli draining.\end{seealsob}

Two tanks exchanging fluid, a drug moving between blood and tissue, and a chain of lakes are all the same object: several reservoirs, one balance equation each, and a linear system. The system machinery of the earlier chapters applies unchanged, and the only genuinely new material is a structural check that catches setup errors before any eigenvalue is computed.

\tech{Compartment models}{hard}
\cue Two or more reservoirs exchanging material: connected tanks, a two-compartment pharmacokinetic model, a lake system, a tracer in blood and tissue.
\method Write one balance equation per compartment, each in the form (in) minus (out), with every flow expressed as a rate times the source compartment's concentration. The result is
\[\mathbf{x}'=A\mathbf{x}+\mathbf{b},\]
solved by the eigenvalue methods for linear systems. The equilibrium is $\mathbf{x}^{*}=-A^{-1}\mathbf{b}$ when $A$ is invertible, and the transient decays at rates set by the eigenvalues of $A$. Before solving, run the conservation check: if the model is closed (nothing enters or leaves), the column sums of $A$ must be zero, $\lambda=0$ must be an eigenvalue, and the corresponding conserved quantity is the total.

\begin{worked}
Two $100$ L tanks are connected so that $10$ L/min flows from the first to the second and $10$ L/min back. Tank 1 starts with $20$ kg of salt, tank 2 with none. With $x=A_1$ and $y=A_2$,
\[x'=\frac{10}{100}y-\frac{10}{100}x=0.1(y-x),\qquad y'=0.1(x-y),\]
so $A=0.1\begin{pmatrix}-1&1\\1&-1\end{pmatrix}$. The column sums are zero, as required for a closed system, and the eigenvalues are $0$ and $-0.2$. Rather than diagonalising, use the two combinations the eigenvalues point to. The sum obeys $(x+y)'=0$, so $x+y\equiv20$. The difference obeys $(x-y)'=-0.2(x-y)$, so $x-y=20\eu^{-0.2t}$. Adding and halving,
\[x(t)=10+10\eu^{-0.2t},\qquad y(t)=10-10\eu^{-0.2t}.\]
At $t=7$ minutes this gives $x=12.4660$, $y=7.5340$, which Runge-Kutta reproduces to six figures. The general lesson is that a symmetric two-compartment system is always solved by sum and difference, and the relaxation rate is $2r/V$, twice what a single-tank intuition would suggest, because both tanks are moving toward each other.
\end{worked}

\begin{worked}[A two-compartment drug model]
A drug is infused into the blood at a constant rate $I$, exchanges with tissue at rates $k_{12}$ and $k_{21}$, and is cleared from the blood at rate $k_e$. Writing $b$ and $q$ for the blood and tissue amounts,
\[b'=I-(k_{12}+k_e)b+k_{21}q,\qquad q'=k_{12}b-k_{21}q.\]
This system is \emph{not} closed, so $\lambda=0$ is not an eigenvalue and there is a genuine steady state. Setting both derivatives to zero, the second equation gives $q^{*}=\frac{k_{12}}{k_{21}}b^{*}$, and substituting into the first collapses the exchange terms entirely: $I=k_eb^{*}$, so $b^{*}=I/k_e$ and $q^{*}=\frac{k_{12}I}{k_{21}k_e}$. The steady blood level depends only on the infusion rate and the clearance rate, not on the exchange rates at all, which is why the tissue compartment changes how fast a drug reaches its plateau but not where the plateau sits.
\end{worked}

\begin{trap}
Omitting the return flow, or writing it with the wrong compartment's concentration. The flow from tank 2 into tank 1 is $r_{21}\cdot A_2/V_2$, using tank 2's concentration, and putting $A_1/V_1$ there produces a matrix whose column sums are not zero. Run the conservation check on every closed model: it costs one addition per column and catches almost every sign and index error in the setup.
\end{trap}

\begin{seealsob}Mixing and tank problems; linear systems and the matrix exponential; Newton's law of cooling.\end{seealsob}

Draining under gravity breaks the pattern of the previous two techniques in one respect: the outflow rate is not given but determined by the physics, and it depends on the state. The result is separable rather than linear, and it produces the chapter's first solution that reaches its terminal state in finite time.

\tech{Torricelli draining and square-root emptying}{medium}
\cue A tank draining through a hole in the bottom, with the outflow speed proportional to $\sqrt{h}$; equivalently any statement that the efflux velocity is $\sqrt{2gh}$.
\method Volume balance with cross-sectional area $A(h)$ at height $h$ and hole area $a$:
\[A(h)\dv{h}{t}=-ca\sqrt{2gh}.\]
For a tank of constant cross-section this is $h'=-k\sqrt h$, separable, and the useful form of the answer is that $\sqrt h$ is \emph{affine in time}:
\[2\sqrt h=2\sqrt{H}-kt,\qquad h(t)=\Bigl(\sqrt H-\tfrac k2t\Bigr)^{2},\]
so the tank empties at the finite time $t_{\text{empty}}=\dfrac{2\sqrt H}{k}$. Plot $\sqrt h$ against $t$ to fit $k$ from data: it is a straight line, and its intercept is the emptying time.

\begin{worked}
A tank of constant cross-section, initially $4$ m deep, obeys $h'=-\tfrac12\sqrt h$. Then $\sqrt h=2-\tfrac14t$, so
\[h(t)=\Bigl(2-\tfrac t4\Bigr)^{2},\qquad t_{\text{empty}}=8\ \text{minutes},\]
and at $t=4$ the depth is $h=1$ m, not $2$ m: half the time drains three quarters of the depth. Numerical integration with the right-hand side cut off at $h=0$ agrees, and confirms that $h$ stays at $0$ after $t=8$.
\end{worked}

\begin{worked}[Non-constant cross-section]
A hemispherical bowl of radius $R$ drains through a hole at the bottom. At depth $h$ the free surface is a circle of radius $\sqrt{R^{2}-(R-h)^{2}}=\sqrt{2Rh-h^{2}}$, so $A(h)=\pi(2Rh-h^{2})$ and
\[\pi(2Rh-h^{2})\dv{h}{t}=-k\sqrt h\ \Longrightarrow\ \frac{\pi}{k}\bigl(2Rh^{1/2}-h^{3/2}\bigr)\dd h=-\dd t.\]
Integrating from $h=R$ down to $0$,
\[t_{\text{empty}}=\frac{\pi}{k}\left[\frac{4R}{3}h^{3/2}-\frac{2}{5}h^{5/2}\right]_{0}^{R}=\frac{\pi}{k}\left(\frac43-\frac25\right)R^{5/2}=\frac{14\pi R^{5/2}}{15k}.\]
The lesson generalises: the geometry enters only through $A(h)$, and every such problem is the single integral $\int A(h)h^{-1/2}\dd h$.
\end{worked}

\begin{trap}
Using the formula past the emptying time. The expression $\bigl(\sqrt H-\tfrac k2t\bigr)^{2}$ is a perfectly good parabola that turns around and increases after $t_{\text{empty}}$, describing a tank that spontaneously refills. The solution of the physical problem is that formula on $[0,t_{\text{empty}}]$ and $h\equiv0$ afterwards, and the two pieces genuinely both solve the equation from the moment they meet: uniqueness fails at $h=0$ because $\sqrt h$ is not Lipschitz there. That failure is the subject of the next chapter, and this is where a reader first meets it in the wild.
\end{trap}

\begin{seealsob}Mixing and tank problems; separable equations; solutions lost by dividing.\end{seealsob}

\section{Mechanical and electrical archetypes}

The models so far accumulated a substance. The next three conserve or balance something else: charge around a loop, momentum against gravity and drag, momentum against an exhaust stream. The mathematics repeats, but the recognition cues are entirely different, and so are the standard errors.

\tech[RC and RL circuits]{RC and RL circuits}{core}
\cue A single-loop circuit with exactly one energy-storage element (a capacitor or an inductor) and a source. Words such as ``time constant'', ``steady state current'', ``transient''.
\method Kirchhoff's voltage law gives a first-order linear equation in each case:
\[L\dv{I}{t}+RI=E(t),\ \ \tau=\frac LR;\qquad R\dv{Q}{t}+\frac QC=E(t),\ \ \tau=RC.\]
For a constant source the answer is always \emph{steady state plus decaying transient}, and you can write it without integrating: identify the final value, identify the initial value, and interpolate exponentially,
\[X(t)=X_\infty+\bigl(X_0-X_\infty\bigr)\eu^{-t/\tau}.\]
For a sinusoidal source use the steady-state method of the constant-coefficient chapter: the amplitude of the current in an RL loop driven by $E_0\sin\omega t$ is $E_0/\sqrt{R^{2}+\omega^{2}L^{2}}$, and the phase lag is $\arctan(\omega L/R)$.

\begin{worked}
$R=12\,\Omega$, $L=4$ H, $E=60$ V, $I(0)=0$. The equation is $4I'+12I=60$, that is $I'+3I=15$, so $\tau=L/R=\tfrac13$ s and $I_\infty=E/R=5$ A. Interpolating,
\[I(t)=5\bigl(1-\eu^{-3t}\bigr).\]
The current reaches $99\%$ of its final value after $t=\tau\ln100\approx1.535$ s, which is the useful engineering statement: five time constants is $99.3\%$, and that is the origin of the rule of thumb.
\end{worked}

\begin{worked}[Discharge, and reading $\tau$ off a graph]
A capacitor charged to $Q_0$ is discharged through a resistor: $RQ'+Q/C=0$, so $Q=Q_0\eu^{-t/RC}$. With $R=100\,\Omega$ and $C=200\,\mu$F, $\tau=RC=100\times2\times10^{-4}=0.02$ s. After $50$ ms the charge is $Q_0\eu^{-2.5}=0.0821Q_0$. Conversely, given a measured decay from $Q_0$ to $Q_0/2$ in $13.9$ ms, $\tau=13.9/\ln2=20.05$ ms, recovering $RC$ without knowing either component. The half-life language of radioactive decay transfers verbatim.
\end{worked}

\begin{trap}
Forming the time constant by the wrong combination. For RC it is the \emph{product} $RC$ and for RL it is the \emph{quotient} $L/R$, and the two are easy to swap under pressure. The unit check settles it in three seconds: ohms times farads is seconds, henries divided by ohms is seconds, and $R/L$ or $C/R$ is neither. The second trap is applying $I=\frac ER(1-\eu^{-t/\tau})$ when $I(0)\ne0$; use the interpolation form, which handles any initial value.
\end{trap}

\begin{seealsob}Mixing and tank problems; Newton's law of cooling; first-order linear equations.\end{seealsob}

A falling body with resistance is the archetype that teaches sign discipline, because the drag force reverses with the motion and the two natural resistance laws produce two entirely different functions. Both are separable; only one is also linear.

\tech{Falling bodies with drag and terminal velocity}{medium}
\cue ``Terminal velocity'', ``air resistance proportional to speed'' or ``to the square of the speed'', a parachute, a raindrop, a body launched upward against resistance.
\method Take down as positive and write $m v'=mg-F_{\text{drag}}$.

\emph{Linear drag}, $F=bv$. The equation $mv'=mg-bv$ is both linear and separable, with terminal velocity $v_\infty=mg/b$ and
\[v(t)=v_\infty\bigl(1-\eu^{-bt/m}\bigr)=v_\infty\bigl(1-\eu^{-gt/v_\infty}\bigr).\]

\emph{Quadratic drag}, $F=bv^{2}$. Now $mv'=mg-bv^{2}$ separates with $v_\infty=\sqrt{mg/b}$ to give
\[v(t)=v_\infty\tanh\!\left(\frac{gt}{v_\infty}\right),\qquad y(t)=\frac{v_\infty^{2}}{g}\ln\cosh\!\left(\frac{gt}{v_\infty}\right).\]
In both cases $v_\infty$ is read off without solving anything, by setting $v'=0$ in the equation. Do that first: it is one line and it is frequently the entire question.

\begin{worked}
$v'=9.8-0.2v$ with $v(0)=0$. Setting $v'=0$ gives $v_\infty=9.8/0.2=49$ m/s immediately, and the time constant is $m/b=1/0.2=5$ s, so
\[v(t)=49\bigl(1-\eu^{-t/5}\bigr),\qquad v(5)=49(1-\eu^{-1})=30.9739\ \text{m/s}.\]
Runge-Kutta integration returns $30.973907$, agreeing to six decimals.
\end{worked}

\begin{worked}[Quadratic drag, and why the approach is faster]
A skydiver with $b/m=0.002$ m$^{-1}$ has $v_\infty=\sqrt{9.8/0.002}=\sqrt{4900}=70$ m/s. Then
\[v(t)=70\tanh(0.14t),\qquad v(10)=70\tanh(1.4)=70\times0.885352=61.97\ \text{m/s},\]
which is $88.5\%$ of terminal after ten seconds. Compare a linear-drag model tuned to the same terminal speed: $v=70(1-\eu^{-gt/70})=70(1-\eu^{-0.14t})$ reaches only $70(1-\eu^{-1.4})=52.75$ m/s, or $75.3\%$. Quadratic drag is weak at low speed and therefore lets the body accelerate closer to terminal before biting, so $\tanh$ approaches its asymptote faster than the exponential does. If a problem gives you the fraction of terminal velocity at a stated time, that fraction alone tells you which drag law was intended.
\end{worked}

\begin{trap}
Keeping the form $-bv^{2}$ while the object is moving \emph{upward}. Drag opposes motion, so the honest force is $-b\,v\abs{v}$, and on the way up, with down still positive, $v<0$ and the drag term is $+bv^{2}$: the equation becomes $mv'=mg+bv^{2}$ on the ascent and $mv'=mg-bv^{2}$ on the descent. These are genuinely different equations with genuinely different solutions ($\tan$ on the way up, $\tanh$ on the way down), and one consequence is that a projectile launched with speed $u$ returns with speed strictly less than $u$. Solve the two phases separately and match at the apex.
\end{trap}

\begin{seealsob}The Tsiolkovsky rocket equation; RC and RL circuits; separable equations.\end{seealsob}

The last mechanical archetype is the one where the naive translation is wrong, and the reason it is wrong is worth more than the formula. When mass leaves the body, $F=ma$ is no longer the correct statement of Newton's law, and the correct starting point is a momentum balance over the body plus the matter it has just ejected.

\tech{The Tsiolkovsky rocket equation}{hard}
\cue A body whose mass changes as it moves: a rocket burning fuel, a chain being picked up off a table, a raindrop accreting vapour.
\method Work with momentum over a short interval, comparing the system before and after the ejection. In free space, with exhaust speed $u$ \emph{relative to the rocket}, conservation of momentum gives $m\dd v=-u\dd m$, which separates immediately:
\[v-v_0=u\ln\frac{m_0}{m}.\]
Under constant gravity and a constant burn rate $\abs{\dot m}$, the same balance gives $m\dv{v}{t}=-u\dot m-mg$, whose solution is the rocket equation minus a gravity loss:
\[v(t)=u\ln\frac{m_0}{m(t)}-gt.\]

\begin{worked}
A single stage with exhaust speed $u=3000$ m/s is to deliver a velocity change of $9000$ m/s. Then $\ln(m_0/m_f)=3$, so $m_0/m_f=\eu^{3}=20.0855$ and the final mass is $4.98\%$ of the initial mass. Since structural mass alone is typically eight to ten percent of a stage, this calculation is the reason single-stage-to-orbit vehicles are hard and staging exists: the logarithm makes velocity expensive in a way no engineering margin can absorb.

The special case worth memorising is $\Delta v=u$, which needs $m_f=m_0/\eu$, so the fuel fraction is
\[1-\frac1{\eu}=0.632121=63.21\%.\]
\end{worked}

\begin{trap}
Starting from $F=\dv{(mv)}{t}$ and expanding it as $m\dot v+\dot mv$. That expression is not frame-independent when mass crosses the system boundary: evaluate it in a frame moving with the rocket and you get a different answer, which is proof that it is not a physical law. The correct relation uses the exhaust velocity relative to the rocket, and the resulting thrust term is $-u\dot m$ with $\dot m<0$, hence positive. Any derivation that never mentions the word ``relative'' has gone wrong.
\end{trap}

\begin{seealsob}Falling bodies with drag; separable equations; logarithmic differentiation.\end{seealsob}

\section{Models that arrive as geometry}

The remaining two archetypes are not physical balances at all. Their statements are geometric conditions on a curve, and the translation step converts a condition about tangent lines into a condition about $y'$. That conversion is the whole difficulty, and the equations that result are usually easy.

\tech{Orthogonal trajectories}{medium}
\cue ``Find the family of curves cutting the given family at right angles'', or a physical pairing that is always orthogonal: isotherms against heat-flow lines, equipotentials against field lines, streamlines against velocity potentials.
\method Three steps, and the middle one is the one people skip.
\begin{enumerate}
\item Differentiate the family $F(x,y,C)=0$ implicitly.
\item \emph{Eliminate $C$} between the original equation and its derivative, producing a slope field $y'=f(x,y)$ in $x$ and $y$ alone.
\item Solve $y'=-1/f(x,y)$.
\end{enumerate}
The elimination is what makes the slope field a property of the geometry rather than of the labelling, and until it is done the negative reciprocal means nothing.

\begin{worked}
The family $y=Cx^{2}$ of parabolas through the origin. Differentiating, $y'=2Cx$; eliminating $C=y/x^{2}$ gives the slope field
\[y'=\frac{2y}{x}.\]
The orthogonal family therefore solves $y'=-\dfrac{x}{2y}$, which separates: $2y\dd y=-x\dd x$, so $y^{2}+\tfrac12x^{2}=k$, a family of ellipses centred at the origin with their major axes along $y=0$. Every one of them meets every parabola at right angles, and the two families together are the standard picture of a dipole-like flow.
\end{worked}

\begin{worked}[When log-differentiation does the elimination for you]
Take $\eu^{y}\tan\!\left(\dfrac x2+\dfrac\pi4\right)=C$. Taking logarithms first turns the product into a sum and makes $C$ vanish under differentiation without any algebra:
\[y+\ln\tan\!\left(\frac x2+\frac\pi4\right)=\ln C\ \Longrightarrow\ y'+\frac{\tfrac12\sec^{2}u}{\tan u}=0,\qquad u=\frac x2+\frac\pi4.\]
Now $\dfrac{\sec^{2}u}{2\tan u}=\dfrac{1}{2\sin u\cos u}=\dfrac{1}{\sin2u}=\dfrac{1}{\sin(x+\tfrac\pi2)}=\sec x$, so the slope field is $y'=-\sec x$ with no $C$ anywhere. The orthogonal family solves $y'=\cos x$, hence
\[y=\sin x+C_1.\]
Numerically differentiating the original family at $x=0.7$ with $C=2$ returns $-1.30729$, and $-\sec(0.7)=-1.30729$: the elimination is verified before the reciprocal is taken.
\end{worked}

\begin{trap}
Taking the negative reciprocal before eliminating $C$. The expression $y'=2Cx$ describes one curve per value of $C$, and $-1/(2Cx)$ is the perpendicular direction along that one curve, not a field on the plane; solving it produces a one-parameter family for each $C$, which is one parameter too many and answers nothing. A second, quieter trap: points where $f=0$ or $f$ is infinite are where the two families have horizontal and vertical tangents, and the orthogonal equation is singular there. Note them rather than dividing through them.
\end{trap}

\begin{seealsob}Pursuit curves and the tractrix; exact equations and integrating factors; implicit general solutions.\end{seealsob}

Pursuit problems are the same kind of translation with a moving constraint instead of a static one. The constraint is always that the tangent to the curve points at something, and turning that sentence into an equation in $y'$ is the entire modelling step.

\tech{Pursuit curves and the tractrix}{hard}
\cue One object always heading directly at another; an object dragged by a taut string of fixed length; a dog chasing along a bank; a missile with pure pursuit guidance.
\method Write the heading condition as a slope, then differentiate to eliminate the pursued object's position.

\emph{Tractrix.} The tangent segment from the curve down to the $x$-axis has constant length $a$. If the curve is $y(x)$ then the horizontal run of that segment is $\sqrt{a^{2}-y^{2}}$, so
\[\dv{y}{x}=-\frac{y}{\sqrt{a^{2}-y^{2}}},\]
which separates. The solution, with the cusp at $(0,a)$, is
\[x=a\ln\frac{a+\sqrt{a^{2}-y^{2}}}{y}-\sqrt{a^{2}-y^{2}}.\]

\emph{True pursuit.} Target moving up the $y$-axis from the origin at speed $v$, pursuer starting at $(a,0)$ with speed $w$, always heading at the target. Writing $k=v/w$ and eliminating time gives
\[x\ddv{y}{x}=k\sqrt{1+\left(\dv{y}{x}\right)^{2}},\]
a second-order equation in which $y$ does not appear, so $p=y'$ reduces it to first order.

\begin{worked}[The tractrix, verified]
Differentiating the closed form with respect to $y$ and writing $s=\sqrt{a^{2}-y^{2}}$,
\[\dv{x}{y}=a\left(\frac{-y/s}{a+s}-\frac1y\right)+\frac ys=\frac{-ay^{2}-as(a+s)+y^{2}(a+s)}{ys(a+s)}=\frac{-s^{2}(a+s)}{ys(a+s)}=-\frac sy,\]
using $y^{2}=a^{2}-s^{2}$ to collapse the numerator. Inverting, $\dv{y}{x}=-y/s$, which is the equation. Numerically, at $y=0.6$ with $a=1$ the closed form gives $\dv{y}{x}=-0.75$ and $-y/\sqrt{1-y^{2}}=-0.6/0.8=-0.75$. The curve meets the $y$-axis at $y=a$ with a vertical cusp and approaches the $x$-axis asymptotically as $y\to0$.
\end{worked}

\begin{worked}[Pursuit, and where the capture happens]
Take the pursuit equation with $p=y'$ and $p(a)=0$, since at $t=0$ the target is at the origin and the pursuer at $(a,0)$ is heading along the $x$-axis. Then
\[x\dv{p}{x}=k\sqrt{1+p^{2}}\ \Longrightarrow\ \frac{\dd p}{\sqrt{1+p^{2}}}=k\frac{\dd x}{x}\ \Longrightarrow\ \operatorname{arcsinh}p=k\ln\frac xa,\]
so $p=\sinh\bigl(k\ln(x/a)\bigr)=\tfrac12\bigl[(x/a)^{k}-(x/a)^{-k}\bigr]$. Integrating once more and imposing $y(a)=0$,
\[y=\frac a2\left[\frac{(x/a)^{k+1}}{k+1}-\frac{(x/a)^{1-k}}{1-k}\right]+\frac{ak}{1-k^{2}}.\]
If $k<1$, so the pursuer is faster, both powers vanish as $x\to0$ and capture occurs on the $y$-axis at
\[y_{\text{capture}}=\frac{ak}{1-k^{2}}.\]
With $k=\tfrac12$ and $a=1$ that is $\tfrac{1/2}{3/4}=\tfrac23$. If $k=1$ the second term integrates to a logarithm instead and $y\to\infty$: the pursuer closes on the target but never catches it, converging to a fixed lag of $a/2$. If $k>1$ the pursuer is slower and never gets there at all.
\end{worked}

\begin{trap}
Setting up the geometry with the wrong sign, so that the pursuer flees. In the tractrix the trailing object is always \emph{behind}, so $y$ decreases as $x$ increases and $y'<0$; a derivation returning $y'=+y/\sqrt{a^{2}-y^{2}}$ has the string pushing rather than pulling. Sketch the configuration and check the sign of $y'$ against the picture before integrating, because both signs integrate perfectly well and only one describes the problem.
\end{trap}

\begin{seealsob}Orthogonal trajectories; reduction of order when the dependent variable is absent; separable equations.\end{seealsob}

\begin{keynote}[The seven shapes, and how to tell them apart in one reading]
Every closed form in this chapter is one of seven functions, and the wording of the problem picks it out before any algebra. \emph{Pure exponential} $N_0\eu^{-\lambda t}$: proportional decay with nothing else happening. \emph{Shifted exponential} $X_\infty+(X_0-X_\infty)\eu^{-t/\tau}$: anything approaching an equilibrium, which covers cooling, RC and RL with a constant source, constant-volume mixing, and decay against a source. \emph{Logistic} $K/(1+A\eu^{-rt})$: growth against a ceiling. \emph{Hyperbolic tangent} $v_\infty\tanh(gt/v_\infty)$: quadratic drag. \emph{Rational or power} $2(50+t)-c(50+t)^{-2}$: mixing with changing volume, where the integrating factor is a power of $V$. \emph{Quadratic reaching zero} $\bigl(\sqrt H-\tfrac k2t\bigr)^{2}$: gravity-driven draining, and the only member of the list that finishes in finite time. \emph{Logarithm} $u\ln(m_0/m)$: variable mass. If your answer is not one of these seven and the problem was not a compartment system, recheck the translation rather than the integration.
\end{keynote}

\begin{keynote}[Steady state first, always]
In every autonomous model here, setting the derivative to zero and solving the resulting algebraic equation costs one line and yields the terminal value: $T_a$ for cooling, $K$ for logistic, $c_{\text{in}}V$ for constant-volume mixing, $E/R$ for the RL current, $\sqrt{mg/b}$ for quadratic drag, $-A^{-1}\mathbf b$ for a compartment system. Do this before solving, for three reasons. It is often the entire question, since ``find the terminal velocity'' and ``find the limiting concentration'' never require the transient. It gives you the substitution $U=X-X_\infty$ that reduces the equation to the homogeneous one. And it is a check on the closed form you eventually produce: if $\lim_{t\to\infty}X(t)$ does not match the equilibrium you computed in one line, the solution is wrong and you know it before the marker does.
\end{keynote}

\section{Recognition matrix}
\begin{recog}{@{}p{0.32\textwidth}p{0.35\textwidth}p{0.27\textwidth}@{}}
\rowcolor{mist}\mh{If the problem says} & \mh{Write} & \mh{Then} \\
\midrule
\endhead
``Half-life'', ``carbon dating'' & $N=N_0 2^{-t/t_{1/2}}$ & Never compute $\lambda$ \\
``Proportional to the amount present'' & $N'=-\lambda N$, pure exponential & Check whether a source term is also present \\
Constant production and decay & $N'=R-\lambda N$; equilibrium $R/\lambda$ & Substitute $U=N-R/\lambda$ \\
Cooling or heating toward surroundings & $U=T-T_a$ decays exponentially & Ratios of \emph{excess}, never of $T$ \\
Three readings, ambient unknown & $T_a=\dfrac{T_0T_2-T_1^{2}}{T_0+T_2-2T_1}$ & Verify the excesses are geometric \\
Two readings at equal spacing & Excess multiplies by a fixed $\rho$ & Answer in powers of $\rho$; skip $k$ \\
``Levels off at a carrying capacity'' & Logistic; odds $\dfrac{P}{K-P}$ grow as $\eu^{rt}$ & $P=\dfrac{K}{1+W_0^{-1}\eu^{-rt}}$ \\
``Fastest growth occurred in year $t_1$'' & That is the inflection: $P(t_1)=K/2$ & Use it as the third data point \\
Well-stirred tank, brine in and out & $A'=r_{\text{in}}c_{\text{in}}-\dfrac{r_{\text{out}}}{V(t)}A$ & Track amount, not concentration \\
Inflow rate $\ne$ outflow rate & $V(t)=V_0+\delta t$; factor $(V_0+\delta t)^{r_{\text{out}}/\delta}$ & Note the time the tank fills or empties \\
Two or more connected reservoirs & One balance equation each; $\mathbf x'=A\mathbf x+\mathbf b$ & Closed system: column sums $0$, $\lambda=0$ \\
Symmetric two-tank exchange & Solve by sum and difference & Relaxation rate $2r/V$ \\
Drug infused, blood and tissue & $b^{*}=I/k_e$, independent of exchange & Exchange sets the speed, not the plateau \\
Draining through a hole & $A(h)h'=-ka\sqrt h$; $\sqrt h$ affine in $t$ & Stop at $t_{\text{empty}}$; then $h\equiv0$ \\
``Time constant'', one storage element & RL: $\tau=L/R$; RC: $\tau=RC$ & Unit check: both must be seconds \\
Step source, any first-order model & $X_\infty+(X_0-X_\infty)\eu^{-t/\tau}$ & No integration needed \\
Sinusoidal source, RL loop & Amplitude $E_0/\sqrt{R^{2}+\omega^{2}L^{2}}$ & Phase lag $\arctan(\omega L/R)$ \\
``Terminal velocity'' & Set $v'=0$ in the equation & One line; often the whole question \\
Resistance proportional to $v$ & $v=v_\infty(1-\eu^{-gt/v_\infty})$, $v_\infty=mg/b$ & Linear \emph{and} separable \\
Resistance proportional to $v^{2}$ & $v=v_\infty\tanh(gt/v_\infty)$, $v_\infty=\sqrt{mg/b}$ & Distance: $\frac{v_\infty^{2}}{g}\ln\cosh$ \\
Body thrown \emph{upward} against drag & Drag is $-b\,v\abs{v}$: two different equations & $\tan$ up, $\tanh$ down; match at apex \\
Mass changes as the body moves & $m\dd v=-u\dd m$, so $v=u\ln(m_0/m)$ & Never expand $\dv{(mv)}{t}$ \\
Rocket under gravity, constant burn & $v=u\ln(m_0/m)-gt$ & The $-gt$ is the gravity loss \\
``Cuts the family at right angles'' & Differentiate, \emph{eliminate $C$}, then $y'\mapsto-1/y'$ & A leftover $C$ voids the answer \\
Family is a product of factors & Log-differentiate: $C$ disappears free & Then simplify before inverting \\
``Always heading directly at'' & Tangent points at the target; then $p=y'$ & Capture at $ak/(1-k^{2})$ if $k<1$ \\
Tangent segment of fixed length $a$ & Tractrix: $y'=-y/\sqrt{a^{2}-y^{2}}$ & Check the sign against a sketch \\
\end{recog}
